\documentclass[11pt,a4paper]{article}

\usepackage[T1]{fontenc}
\usepackage[utf8]{inputenc}
\usepackage[spanish]{babel}
\usepackage[a4paper,margin=2.8cm]{geometry}
\usepackage{microtype}
\usepackage{xurl}
\usepackage[hidelinks]{hyperref}

\hypersetup{
  pdftitle={Neutrinos, mentiras y detectores de IA. No hay atajos a Ítaca},
  pdfauthor={Juan José Gómez Cadenas}
}

\setlength{\parindent}{1.25em}
\setlength{\parskip}{0.35em}
\raggedbottom

\title{Neutrinos, mentiras y detectores de IA\\[0.4em]\large No hay atajos a Ítaca}
\author{Juan José Gómez Cadenas}
\date{}

\begin{document}

\maketitle

Lo primero que se me viene a la cabeza, cuando me entero del premio Nobel a Francis Halzen, es la entrevista que Francesc y yo le hicimos en 2014\footnote{\url{https://www.jotdown.es/2014/05/francis-halzen-siempre-aconsejo-a-mis-estudiantes-no-leais-demasiados-libros-haced-cosas/}} en la que Francis, siempre un iconoclasta, aconsejaba a sus alumnos no leer demasiados libros.

Curioso consejo viniendo de un científico de primera clase, ¿no? Y sin embargo Francis sabía lo que se decía. El gran detector enterrado bajo la Antártida, IceCube\footnote{\url{https://www.jotdown.es/2012/12/juan-jose-gomez-cadenas-paisaje-con-neutrinos/}}, funciona porque los físicos que lo construyeron ignoraron las bien documentadas medidas que aseguraban que la luz apenas se propaga unas decenas de metros en hielo. IceCube reparte sus más de cinco mil fotomultiplicadores ---86 cadenas con 60 detectores ópticos gigantes cada una--- en un volumen de un kilómetro cúbico aproximadamente. La distancia promedio entre esas cadenas es casi de cien metros. Según los libros, la señal de una posible interacción de neutrinos cósmicos no podría detectarse en el aparato porque la luz que liberaría en el hielo no alcanzaría esos fotomultiplicadores y no produciría por tanto una señal observable.

Pero resulta que la atenuación de la luz en el hielo ultrapuro que se encuentra a un kilómetro por debajo del nivel del suelo en la Antártida es mucho menor que en el hielo de la superficie, mucho menor que en cualquier hielo que se conocía, mucho menor de lo que estaba escrito en los libros. Y así, gracias a que Francis y sus colegas ignoraron un dato que no aplicaba a su problema ---no por arrogancia, naturalmente, sino por intuición al principio y por refinadas medidas más tarde--- gracias al arrojo de un físico teórico, que «no sabía» ---es decir, no quería saber--- que las dificultades logísticas de construir semejante aparato en el Polo Sur eran inabordables, se construyó el detector más ambicioso y bello de la historia ---aunque reconozco que las bestias que dormitan en el túnel del LHC no son tampoco despreciables dragones---. Un detector que ha visto con sus ojos electrónicos la explosión de luz causada por neutrinos de energías fabulosas, viajeros galácticos e intergalácticos que nos hablan de las catástrofes inimaginables que los produjeron.

Nadie en IceCube ---ni en mi experimento, NEXT\footnote{\url{https://next-experiment.org/}}, que aspira a descubrir algo tan insólito como los neutrinos cósmicos, a saber, que el neutrino es su propia antipartícula--- está muy preocupado, que yo sepa, por la posibilidad de que la IA nos quite el trabajo. Naturalmente, todos usamos redes neuronales para clasificar los sucesos que registramos en nuestros aparatos, de hecho, algunos las usamos desde hace más de dos décadas, desde los tiempos heroicos donde manejar esos algoritmos levantaba rumores y sospechas. Por supuesto, todos mis estudiantes de doctorado y post-docs ---incluyéndome a mí mismo--- tienen a Claude o Codex de fieles compañeros de programación. Y por supuesto, el uso de estas herramientas supone, como cualquier avance tecnológico, un aliciente y un reto. Mis estudiantes tienen ahora un mentor que les ayuda a entender con detalle los problemas en los que trabajan ---quién lo hubiera pillado a su edad--- pero también corren el riesgo de volverse perezosos y dejar que sea Claude quien haga las cuentas. Los más viejos vemos cómo nuestra experiencia e intuición pueden reforzarse con este esclavo digital ---yo prefiero llamarlo socio y tratarlo de colega, en parte porque los buenos modales no hacen daño a nadie, en parte por comprar boletos para el indulto después de la revolución de Skynet--- tan simpático que hace todo lo que le pedimos, eso sí, a riesgo de volvernos un poco majaras. 

Son tiempos interesantes para los físicos experimentales, pero la IA no le sugirió a Francis construir su detector bajo la Antártida ---de hecho, si la IA hubiera funcionado por aquellas fechas le habría aconsejado lo contrario--- ni nos explicó a James White, Dave Nygren y servidor que se podía construir un detector para buscar desintegraciones doble beta sin neutrinos con gas xenón a alta presión ---en nuestro caso, posiblemente nos habría desanimado alegando, con toda razón, que la tecnología estaba inmadura y aconsejándonos probar con xenón líquido---. La IA ---es decir, los grandes modelos de lenguaje o LLM que llamamos IA, confundiéndolos con otras aplicaciones menos sabelotodo y sobre todo nada charlatanas, como las redes convolucionales que usamos para identificar la topología de nuestra señal en NEXT--- es muy competente cuando se las tiene que ver con problemas del tipo \emph{desconocidos conocidos}. Si algo está en la literatura, o se puede conectar ---a menudo de manera no obvia--- a partir de lo que hay en la literatura, o admite deducción lógica, la IA suele tener mucho que decir. Cuando se trata de \emph{desconocidos desconocidos}, es decir, problemas que no están documentados, o, más interesante aún, problemas que ni siquiera están formulados, los LLM pueden servir de ayuda, como la mezcla un poco impía de calculadora y enciclopedia hinchada de esteroides que son, pero dudo que le quiten el trabajo a los humanos.

Cuando P.A.M. Dirac se da cuenta de que la ecuación que ha escrito para describir el movimiento relativista de los electrones contiene soluciones de energía negativa, propone una salida al dilema que a sus colegas ---incluyendo a Heisenberg, uno de los padres de la Mecánica Cuántica--- les parece una memez y que a Claude o a Codex ---de haber existido por la época, alimentados por los artículos de Heisenberg, Pauli, De Broglie, Einstein y otros genios, además de los ejércitos de mediocres que constituyen el grueso de los batallones de la Ciencia--- les habría parecido igualmente tonta. Dirac postuló nada menos que la existencia de la antimateria, un concepto poco menos que blasfemo. ¿Por qué demonios se habría molestado la naturaleza en construir un espejo de sí misma, creando para cada partícula una copia idéntica, con todas las cargas invertidas, una anti-Alicia al otro lado de un espejo fatal listo para aniquilarse con su Alicia?

Y sin embargo, la antimateria existe. Y el fantasma llamado neutrino que Pauli inventó como remedio desesperado para salvar la conservación de la energía ---un invento que despertó no pocas suspicacias y del que el propio Pauli llegó a lamentarse--- también resultó existir, aunque los LLM lo hubieran negado más que Pedro a Cristo. Y algunos, como el que suscribe, estamos seguros ---y apostamos nuestras carreras en ello--- de que el neutrino es su propia antipartícula, como postuló el malogrado Ettore Majorana\footnote{\url{https://www.jotdown.es/2013/02/juan-jose-gomez-cadenas-paisaje-sin-neutrinos-i-un-remedio-desesperado/}}. Podría escribir muchas páginas hablando de los desconocidos desconocidos, que van desde el descubrimiento de la radiación de fondo de microondas hasta la penicilina, pero estos ejemplos que cito son particularmente relevantes porque no son revelaciones accidentales de la naturaleza, sino saltos increíbles de la mente humana. Ahí están también Einstein, Gödel, Galileo y Newton, por no hablar de Demócrito y Epicuro y por ceñirme solo al campo de la física ---con la licencia de Gödel, que me da muchos años de admirarlo---. ¿Llegará algún día la Superinteligencia en ciernes a suplantar esos saltos de la intuición humana? No hay forma de saberlo, claro está, nadie puede leer el futuro. Pero dudo que los actuales LLM lleguen a ello, ni de lejos. Los LLM se alimentan de texto y los desconocidos desconocidos no han producido aún texto que permita meterlos en vereda.

Y sin embargo, me manda mi hijo noticia de que Hugo Duminil-Copin, medalla Fields de 2022, confiesa haberse despertado llorando, desesperado, tras el reciente y famoso «asalto» de OpenAI a las matemáticas ---o más bien a los matemáticos---. Su caso no es único. La mitad del gremio está de luto, la otra mitad en ese estado de negación que sigue a un funeral. Mi mujer, física teórica de gran prestigio, está deprimida igual que todos mis amigos de ese gremio. Todos se ahogan en la insoportable levedad del ser en que los ha sumido la IA.

Se entiende, claro está. Mi trabajo diario ---que es el que cuenta--- depende de entender un detector único en el mundo y para eso no hay IA que valga. No son pocas las veces en las que, hechos un lío con alguna de las rarezas de NEXT ---la velocidad de deriva variable, su mala costumbre de chispear cuando no debe--- nos hemos embarcado en sesudas consultas con Claude, que ha acabado por recomendarnos que leamos... ¡nuestros propios artículos! De donde no hay, no se puede sacar y la experiencia de un equipo como el nuestro, tras quince años de sudar tinta china en el subterráneo, no está contenida en ninguna matriz de pesos. Y eso que la IA se conoce todos nuestros artículos de memoria. Pero una cosa es predicar y otra entender las chispas del cátodo a cincuenta mil voltios.

En cambio, el trabajo diario de mi santa y sus colegas pasa por laboriosos cálculos y esos cálculos, que viven en el mismo territorio que las demostraciones matemáticas de Hugo Duminil-Copin, están al alcance de los LLM actuales. No creo que haya nada de místico en ello. Por el contrario, pienso que la explicación se halla en que muchos de los problemas de física teórica y matemáticas son desconocidos conocidos. No sabemos resolverlos, pueden ser laboriosos, quizás se hayan resistido durante décadas al asalto de los científicos, pero hay suficiente bibliografía, suficientes ataques documentados al problema, suficientes cálculos periféricos y aproximaciones tangenciales relacionados con este, suficiente conocimiento, en resumen, codificado en los pesos de los gigantes modelos ---aún más gigantes si se trata de las versiones privadas que los tecno-oligarcas usan y escamotean al común de los mortales--- como para permitir que un ataque que combine la Enciclopedia de Babel con el músculo de Terminator pueda resolverlo. En mi opinión, los famosos setecientos problemas resueltos por IA que ha publicado recientemente OpenAI\footnote{\url{https://openai.com/index/sharing-ai-progress-in-mathematics/}} y han dejado KO a los matemáticos, o los cálculos a tres bucles cuánticos que la IA se merienda como quien come pipas, son ejemplos espectaculares de desconocidos conocidos, donde la existencia de un gigantesco corpus teórico y la fuerza bruta con la que cuentan los tecno-oligarcas garantizan muchas posibilidades de éxito a los modelos.

De ahí se derivan muchas preguntas válidas y creo que costará un tiempo empezar a entender las respuestas. Hay una parte de orgullo herido, dura y pura. Como ocurrió antes con los ajedrecistas, a los matemáticos se les han metido las IA de okupas y los han echado del piso, o relegado al trastero. Semejante invasión obliga a plantearse aquello de quiénes somos, a dónde vamos y ---como siempre dicen los vascos--- qué comemos hoy. La respuesta del mundo del ajedrez fue bastante contundente. Después de una temporada de duelo, la gente se adaptó sin más. Ciertas cosas han cambiado. Si uno quiere jugar al «ajedrez de verdad» (lento), se va al club, porque en Internet, jugar partidas más largas de 5 minutos es arriesgarse a que le tomen a uno el pelo los tramposos que usan motores (ellos sabrán por qué). Pero no todo es negativo. La IA ha permitido la generación más precoz y brillante de ajedrecistas de la historia. El ajedrez humano goza hoy de mejor salud que nunca, básicamente porque los ajedrecistas han decidido usar las máquinas para entrenar ---es decir, como una herramienta---, y, una vez asimilado que juegan mejor que nosotros, «pasan de ellas».

Los matemáticos tienen muchas opciones posibles. Una es la del propio Duminil-Copin, que ha decidido, por su parte, pasar de los modelos\footnote{\url{https://www.unige.ch/~duminil/publi.html}}. En su página web dice:

\begin{quotation}
He decidido no recurrir a la inteligencia artificial como fuente de ideas nuevas en mi propia investigación. El espectacular progreso de la inteligencia artificial abre oportunidades sin precedentes para ampliar el alcance y las aplicaciones de nuestra disciplina. Sin embargo, ya sea intercambiando ideas con otros matemáticos, enseñando, dirigiendo a estudiantes o compartiendo las matemáticas con el gran público, lo que desempeña el papel principal no es la respuesta en sí, sino el camino que conduce a ella. Por eso deseo seguir siendo, en cierto sentido, un matemático artesano, que se toma el tiempo de vagar, junto a sus colegas, por los rincones ocultos del paisaje matemático.
\end{quotation}

Otros, como Terence Tao, proponen mirar las cosas con cierto contexto. Para él, la IA es una herramienta más, la última ---eso sí, a una escala sin precedentes--- de una larga cadena de herramientas. El influyente matemático ---igual que muchos científicos entre los que me incluyo--- no está en contra de los LLM ni tampoco cree que sus recientes proezas impliquen que la singularidad superinteligente esté a la vuelta de la esquina. En una charla reciente\footnote{\url{https://teorth.github.io/tao-web/slides/math-2-0-caltech-2026.pdf}}, apunta, con razón, que las matemáticas se caracterizan por varios elementos que las hacen una «presa» sabrosa para la «depredación IA». Estas características incluyen: 1) el hecho de que las matemáticas puedan verificarse objetivamente, en particular usando sistemas automáticos como el célebre \emph{Lean}; 2) el hecho de que las matemáticas son un campo casi totalmente digitalizable. En los viejos tiempos, papel y lápiz era todo lo que se necesitaba. Hoy todo el mundo trabaja con ordenadores, pero el concepto no ha cambiado. No hay fugas de gas, ni fusibles que saltan ni aviones que se caen si no funciona el interruptor de un alerón. No hacen falta experimentos (físicos) para verificar sus resultados. Lo que sí hace falta es mucha bibliografía, porque a menudo la solución de un problema se apoya en haber asimilado todo lo que se ha hecho antes en esa área. Es decir, las matemáticas son el dominio perfecto para los LLM y los tecno-oligarcas, que lo saben, han puesto todos los recursos necesarios para que los modelos entren a degüello en el otrora palacio de cristal de las matemáticas. 

Apunta Tao también que otras disciplinas ---sin ir más lejos la física subterránea a la que me dedico--- son menos susceptibles de invasión, al menos hasta que las IA no hayan aprendido a conectar bien un cable o apretar una junta de metal como toca. De hecho, cuanto más se aleja el problema del dominio teórico y más se mete en el mundo, más limitados ---que no inútiles--- son los LLM. Ese efecto lo vemos en ajedrez. Las IA que los ajedrecistas llamamos «motores» solo sirven para jugar al ajedrez y por lo demás son unos muermos ---de ahí que los ajedrecistas usemos un término que revela a las claras su naturaleza mecánica---. Los LLM, por su parte, que según la prensa más calenturienta ya están cerca de alcanzar la superinteligencia, no solo no juegan igual de bien que un motor cualquiera, es que incluso un ajedrecista del montón como servidor les da para el pelo. Este fenómeno, que se suele describir diciendo que el «paisaje de la inteligencia IA tiene picos y valles», permite, por otra parte, poner los pies en el suelo. Si un LLM no da pie con bola jugando al ajedrez, su inteligencia «general» es bastante limitada, por mucho que resuelva problemas del milenio. 

¿Qué problema hay entonces? El mismo Tao apunta que las matemáticas son infinitas y ningún modelo las va a agotar. ¿Por qué no, simplemente, «pasar de ellos», o usarlos como en ajedrez, como en programación, como en tantas otras aplicaciones como las poderosas herramientas que son?

Parte del problema es social. Los tecno-oligarcas se nutren del hype y la prensa les hace el juego de maneras que, a veces, dan vergüenza ajena. Incluso al más calmado de los matemáticos, si le bombardean todo el día con el fin del mundo, acaba por deprimirse. Cabrea también, y no poco, que los señoritos de Silicon Valley, después de apropiarse de toda la ciencia universal ---sin pedir permiso--- se guarden sus modelos pata negra para ellos solos. Pero, hasta cierto punto, se trata de problemas menores. Hay otros más serios y basta recordar los versos de Cavafis para entenderlo. Copio aquí la versión de Lluís Llach, cuyo LP de 1975 se quedó en mi memoria ---y en mi alma--- para siempre:

\begin{verse}
Cuando salgas de viaje hacia Ítaca,\\
has de pedir que el camino sea largo,\\
lleno de aventuras, de conocimiento.\\
Has de pedir que el camino sea largo,\\
que sean muchas las madrugadas,\\
que entrarás en un puerto que tus ojos ignoraban\\
y visitar ciudades para aprender de los que saben.
\end{verse}

Probar un teorema, asegura Tao, es un viaje a Ítaca. Llegar a ella es importante, pero quizás no tanto. Mucho más es el conocimiento, la sabiduría, la experiencia ---y el placer de la belleza--- que se adquiere en el camino. Y el problema de que el LLM nos adelante por la izquierda y llegue tomando un atajo es que rompe las reglas del juego, cambia a los marinos por contables y a los poetas por locutores de reality shows. Como diría el propio Llach, \emph{no és això, companys, no és això}.

La Asociación para las Matemáticas Humanas\footnote{\url{https://www.ahmath.org/}} lo pone negro sobre blanco:

\begin{quotation}
Las herramientas de inteligencia artificial suponen una amenaza fundamental para la práctica científica y la comprensión matemática. El progreso matemático requiere pericia humana, cultivada mediante la escritura y el pensamiento cuidadoso. La producción y el consumo automatizados de matemáticas socavan, por tanto, los cimientos intelectuales de nuestro campo. Los objetivos de las empresas que desarrollan estas tecnologías no están alineados con los de la comunidad matemática; esta industria ha elegido las matemáticas para publicitar las capacidades de sus productos, sin importarle las consecuencias para nuestra disciplina. Los matemáticos tienen, por tanto, la obligación de actuar.
\end{quotation}

También les preocupa (como a mí) el impacto en la educación y la avalancha de basura generada por IA que está inundando las revistas científicas.

\begin{quotation}
Las instituciones matemáticas son especialmente vulnerables a la disrupción. La introducción de herramientas de IA pone en peligro la formación de grado y de posgrado. Las revistas se enfrentan a una explosión de envíos coescritos con IA, y los procesos de publicación están en plena transformación. Los estudiantes están perdiendo oportunidades de adquirir la comprensión profunda que se espera de un matemático. La AHM está organizando varios grupos de trabajo para elaborar propuestas de políticas y herramientas con las que combatir estos problemas.
\end{quotation}

La disrupción que la invasión de la IA supone para la ciencia es real y no se puede ignorar. Quizás en una década parte de estas ansiedades se hayan disipado, como ha ocurrido en el mundo del ajedrez, pero entre tanto hay problemas reales que es necesario atajar.

Uno de ellos, que generaliza el culebrón que están viviendo los matemáticos, es la publicación de resultados científicos. Aquí el problema no son setecientas demostraciones caídas del cielo en un repositorio de GitHub ---OpenAI desdeñando el proceso de publicación formal, posiblemente porque son demasiado guapos para seguir las mismas reglas que todo el mundo--- sino lo contrario. El problema es la avalancha de artículos cuyos autores (o coautores) son en muchos casos «intrusos» con poca o ninguna experiencia en el campo, cuyo trabajo se ha realizado exclusivamente a golpe de prompt ---incluyendo escribir el artículo---.

¿Y por qué eso es un problema? Los más radicales podrían contestarme que cada uno es libre de usar la IA como quiera y que si los resultados son buenos, da igual quién sea el autor o coautor, da igual que lo haya escrito un humano o una máquina.

El argumento vale en teoría y fracasa en la práctica. Veamos por qué con un ejemplo negativo. Este año no espero más de, digamos, 20 artículos en el mundo, en total, dedicados a \emph{resultados}, en la búsqueda de desintegraciones doble beta sin neutrinos. ¿Por qué? Porque solo hay cuatro o cinco experimentos en el mundo en ello y cada uno de nosotros publicamos unos cuatro o cinco artículos al año. Cada uno de esos artículos suele implicar un año de toma de datos, meses de análisis y meses de escritura cuidadosa, revisada por toda la colaboración. Lo mismo puede decirse de IceCube. El detector más prodigioso de la historia de la física, cuyo líder acaba de ganar un premio Nobel, publica unas pocas decenas de artículos al año. El resultado es doble: como el volumen es pequeño, una revista científica tiene poco problema en gestionarlo y como los artículos son de alta calidad ---debido al largo proceso de elaboración--- los árbitros (referees) que tienen que revisarlos pueden hacerlo en unas pocas horas. El proceso de revisión es, normalmente, bastante fluido. El referee pregunta las cosas que no le quedan claras, los autores, que saben de qué hablan y cuentan con el apoyo de toda la colaboración, contestan rápido y el trabajo se publica en un tiempo razonable.

Consideremos ahora artículos sobre \emph{modelos} que predicen diferentes escenarios físicos para que se produzca desintegración doble beta sin neutrinos (o que especulan sobre fuentes de neutrinos cósmicos). Esa es la provincia de los físicos teóricos, que tienden a ser más prolíficos que nosotros, porque pergeñar un modelo ---o, en casi todos los casos, una variante de un modelo existente--- es más barato en términos humanos y técnicos que construir un detector bajo el hielo de la Antártida o en el laboratorio subterráneo de Canfranc. Así que en lugar de 20 artículos, los teóricos quizás produzcan cien en un año. Pero se trata de un negocio todavía manejable, porque en general los autores son científicos con una sólida trayectoria profesional, formada a lo largo de décadas, y una reputación que defender. Pocos de ellos quieren exponerse al ridículo de publicar un artículo incorrecto o banal. 

¿Pero qué ocurre si a la revista no llegan cien, sino mil, o diez mil artículos, la mayoría de ellos de desconocidos, la mayoría de ellos escritos por IA? Ocurre, en primer lugar, que no hay bastantes referees. En segundo lugar, que muchos de esos artículos son irrelevantes, contienen errores o muestran la confusión de quien se ha metido en camisa de once varas armado con Codex y poco sentido común. En tercer lugar, que los referees acaban hartos de tanta basura y acaban por no aceptar un trabajo que todos debemos hacer... o, peor aún, por usar IA para analizar IA. Lo cual no es una solución lógica, como podría suponerse, sino una estupidez al cuadrado.

Solo hay una manera de cortar el slop y quizás alguno me pueda acusar de ponerle puertas al campo y otro de policía de la IA. Pero creo que es imperativo que se usen filtros, para evitar que el sistema reviente. Y el más sencillo de esos filtros es exigir que el artículo pase un detector de IA. Mis argumentos son los siguientes: 

Primer argumento. Frente a la idea de que «si el resultado es bueno, da igual quién lo escriba», la respuesta es: para establecer si el resultado es bueno hace falta el esfuerzo denodado de uno o más humanos (los referees) que no pueden delegar en otra IA, so pena de: a) propagar la ignorancia ---por ejemplo, la IA autora puede decidir por error que la fusión fría es factible y la IA revisora, con la misma matriz de pesos y el mismo error en su condicionamiento, le da la razón--- y b) aumentar la confusión ---la IA revisora estará programada por el referee con el prompt de encontrar pegas y se las encontrará al artículo por perfecto que sea, escribiendo un informe de quinientas páginas llenas de chorradas pedantes y observaciones obvias---. 

Pero ese esfuerzo, al ser un bien escaso, tiene que ser selectivo. Es imposible leerse diez mil artículos y si 9900 de ellos están escritos con IA, ese es un excelente criterio para rechazarlos. Después de todo, a un artículo también se le pide claridad, responsabilidad y \emph{respeto}. Si un autor no respeta lo suficiente a los referees, editores y colegas como para tomarse la molestia de escribir él mismo su artículo, ¿por qué debería un referee, editor, o colega tomarse la molestia de leerlo? Así que si el resultado es bueno, se le pide al autor que haga el esfuerzo de escribirlo él mismo, con lo que de paso, igual se entera de lo que manda a publicar. 

¿Y si su inglés no es bueno? Pues haber estudiao, hombre. Hoy en día no se puede pretender ser científico y no manejarse con el inglés escueto y sin adornos que una buena publicación requiere. Si quiere usar la IA para que le corrija el spelling o los idioms del inglés, que se las traen, que lo haga, ningún detector de IA se va a enfadar por eso. Pero una cosa es el spelling y otra que te escriba el texto la IA. «Es que se organiza mejor»: haber estudiao. «Es que lo explica todo más claramente»: haber estudiao. En ajedrez, a nadie se le ocurre ir a un torneo con su motor en el iPhone diciendo que el motor va a jugar por él porque tiene más ELO.

Segundo argumento, frente a la variante: «las ideas son mías, solo he usado la IA para escribir». No me fastidie, compañero. Si alguien ha invertido las decenas (cientos) de horas que un trabajo serio en física necesita, puede invertir un diez por ciento extra en escribirlo. No hacerlo es como empeñarse en llegar a Ítaca por un atajo imposible. El resto es tomar a tus colegas por tontos.

Así que creo que la policía IA no solo es necesaria, sino esencial, para evitar que los bárbaros invadan el ágora, o la invadan todavía más de lo que ya está. Lo que me lleva a los detectores de IA, que son, hasta cierto punto, detectores de mentiras. Hay una reacción natural en contra de su uso, que creo que está justificada en ciertas situaciones, pero no cuando se trata de analizar publicaciones enviadas a revistas científicas. Los mejores de estos detectores están optimizados, 
como los médicos para «non nocere», no hacer daño. Si se tienen que equivocar, se equivocan dando humano por máquina, porque su obsesión es huir de los falsos positivos. Y hay al menos uno de ellos, Pangram, que documenta, vía un estudio externo, muy pocos falsos positivos\footnote{\url{https://link.springer.com/article/10.1007/s40979-026-00226-w}}.

Y eso, de hecho, es una condición \emph{sine qua non}. Si una revista científica usa Pangram para controlar si los artículos han sido escritos o no por IA, puede poner un límite en, digamos, 50~\% IA (yo personalmente estoy a favor de ponerlo en 10~\% IA como absoluto máximo) y eso solo rechaza casi el 100~\% de los artículos generados por máquina. Si alguno de estos artículos, o porque el autor se ha molestado en reescribirlo mucho, o por carambola, cuela, pues aquí paz y allá gloria. El grueso del ataque ha sido contenido, y eso es lo que cuenta.

Por otra parte, este artículo, como tantos otros míos, se publica en Jot Down. Cabe preguntarse si debería utilizarse el mismo criterio para filtrar las contribuciones de los autores de la revista, como de hecho se está haciendo en Substack. 
Aquí debo reconocer que mi razón y mis tripas no coinciden. En general, en cuanto detecto IA en un artículo dejo de leer. Pero lo cierto es que hay gente que maneja bien la IA y es capaz de usarla para multiplicar y mejorar contenidos, igual que yo la uso para mejorar nuestras investigaciones. Y hay casos en los que el autor, IA o no, consigue que el tema sea verdaderamente interesante, como este\footnote{\url{https://harmoniousdiscourse.substack.com/p/who-owns-discovery}}, que leí ayer, precisamente sobre el problema de la IA y las matemáticas ---cien por cien IA según Pangram y aun así, claro e interesante---.

Así que mi razón me dice que en este caso no hay motivo que justifique exigir que el texto pase un detector de IA. No porque piense que no funcionan ---ya he argumentado que estoy bastante convencido de que si se equivocan se equivocan hacia el lado de no hacer daño--- sino porque aquí no es necesario cortar ningún incendio. El editor general de Jot Down no necesita preocuparse de cómo gestionar diez mil artículos, porque solo va a aceptar ---no le queda otra--- unos pocos al día, y a la hora de decidir quién y qué publica va a seguir su criterio ---que más le vale que sea bueno porque le va en ello el negocio---. Lo mismo puede decirse del autor que decide usar la IA. Si se le va la mano y sus artículos empiezan a ser un poco plastas, sus lectores le darán plantón, así que si tira de Claude es ---la expresión me encanta--- \emph{à ses risques et périls}.


Entre tanto, los neutrinos que ha escupido el agujero negro de la galaxia del Calamar ---hace cuarenta y siete millones de años, aunque nos llegan ahora--- siguen chispeando en el hielo de IceCube, donde según los libros y los LLM no deberían dejar señal. Quizás debemos contemplar la magnitud de lo que significan esas interacciones. Significan un vórtice de estrellas devoradas sin cesar, estrellas en las que quizás hubo planetas donde alguna vez habitaron animales inteligentes, sensibles y dementes como la tribu de monos locos que ha inventado la IA en el nuestro. Cada neutrino que estalla en el Polo Sur nos cuenta una tragedia inenarrable, en tinta azul Cherenkov. O, como escribía mi yo de hace ya más de diez años, sobre IceCube:

\begin{quotation}
Y así, el físico se conmueve, contemplando su paisaje con neutrinos y ruega para que Dios siga mirando a otra parte, lejos de su verde planeta sobrecalentado, lejos de su querida familia a la que no dedica bastante tiempo, lejos de su amado país, destrozado por una raza vil de hombres tenaces, con carnets en los bolsillos que deletrean siglas en las que ya nadie cree. Lejos, lejos, reza el físico, consciente de que no se puede esperar otra cosa de la divinidad que la indiferencia o la aniquilación.

{\itshape Yes, the newspapers were right: snow was general all over Ireland. It was falling softly upon the Bog of Allen and, further westwards, softly falling into the dark mutinous Shannon waves. It was falling too upon every part of the lonely churchyard where Michael Furey lay buried. It lay thickly drifted on the crooked crosses and headstones, on the spears of the little gate, on the barren thorns. His soul swooned slowly as he heard the snow falling faintly through the universe and faintly falling, like the descent of their last end, upon all the living and the dead.}
\end{quotation}

También los neutrinos atraviesan suavemente el universo, también llueven, indiferentes, sobre todos nosotros, los vivos y los muertos.

\end{document}
